\chapter{Weekly calendar, working-hours tab, bulk operations}
\label{ch:calendario-ore-en}

\lettrine[lines=3]{T}{he most}\ substantial UI rewrite since the
project's inception happened in April-May MMXXVI. Three threads
of work converged into a single editorial experience: the
\emph{weekly calendar} with drag-and-drop, the \emph{working-hours
tab} for visual slot configuration, and the \emph{bulk
operations} on the assignments page. The chapter walks them
through in the order in which a user encounters them, starting
from data and ending at the grid.

\section{The \texttt{WeeklyCalendarView} component}

\subsection{Origin and reach}

Until commit \texttt{c8075c0} the \texttt{/schedule} page showed
a static tabular matrix: rows = hours, columns = days, cells
filled with a short ``\textsc{subject} $\cdot$ teacher'' string.
Editing happened only through modals or different tabs: no
direct gesture, no drag. The new interface flips the paradigm:
a \emph{single} Svelte component, \texttt{WeeklyCalendarView},
acts as the shared grid for every calendar view of the product,
in modes \texttt{view}, \texttt{edit} and \texttt{schedule}.

\begin{defbox}[title={The \texttt{WeeklyCalendarView.svelte} file}]
A single source of truth for the weekly grid. It reads days
and slots from \texttt{workingHoursStore} (live propagation,
see \S\ref{sec:store-en}), tunes itself to the
\texttt{mode} prop and changes behavior accordingly. Sixteen
hundred lines of Svelte, but a single hierarchy: \emph{slot}
${\rightarrow}$ \emph{event} ${\rightarrow}$ \emph{action}.
\end{defbox}

\subsection{Three modes}

\begin{description}
  \item[\textsc{view}] Read-only consultation calendar by
    teacher, class or room. No drag, no actions.
  \item[\textsc{edit}] (working-hours tab). The grid does not
    represent lessons but \emph{available slots}. The user
    draws a slot by dragging the mouse top-down
    (\texttt{grab}/\texttt{grabbing} cursor), resizes by
    grabbing the bottom edge (\texttt{ns-resize} cursor), moves
    by grabbing the body, deletes with a click.
  \item[\textsc{schedule}] (\texttt{/schedule}). The grid
    contains the actual lessons of the solved problem. Four
    gestures are allowed --- \emph{Edit}, \emph{Move},
    \emph{Unbind}, \emph{Delete} --- and drag-drop covers both
    the relocation of an already-scheduled lesson and the
    placement of a \emph{not-scheduled} lesson taken from the
    pool sidebar.
\end{description}

\subsection{Anatomy of a slot and conflict preview}

Each grid cell is a \texttt{<div>} with class \texttt{cal-slot}
and \texttt{data-testid="sched-slot-\{D\}-\{H\}"}, where
\texttt{D} is the day index (0$=$Mon, $\dots$, 5$=$Sat) and
\texttt{H} the hour index. Lesson cards live on top of the slot
(\texttt{cal-event} class, \texttt{data-testid="sched-lesson-\{id\}"});
each carries three micro-actions in the top-right corner: pencil
(edit), arrow (unbind), trash (delete). Three distinct CSS
cursors accompany the three drag states (\texttt{grab} at rest,
\texttt{grabbing} during drag, \texttt{ns-resize} on the bottom
edge in edit mode), as dictated by commit \texttt{d4bc873}.

When the user grabs a lesson and hovers the grid
(\texttt{dragover}), \texttt{WeeklyCalendarView} computes in
real time, across the whole weekly matrix, the set of cells that
would generate a \emph{hard conflict} with the dragged lesson
(same teacher in two places, same class in two places, same
classroom occupied). Risky cells take a faint red wash (CSS
class \texttt{cal-slot--drop-conflict}); the source event shows
a small \emph{!}\ red badge
(\texttt{cal-event-conflict-badge}). This is the \emph{soft
conflict preview} introduced by commit \texttt{876755b}: the
user sees what the action would do \emph{before} doing it.

\subsection{The unscheduled-lesson pool}

To the right of the grid, in \texttt{schedule} mode, a sidebar
(testid \texttt{schedule-pool}) lists the lessons that exist in
the database but have no \texttt{day\_idx}/\texttt{hour\_idx}:
typically lessons generated by the optimizer in suspended state
or lessons \emph{unbound} by the user via the homonymous action.
Each pool item is draggable and can be dropped on any slot to
schedule it immediately.

\subsection{The four schedule actions}

\begin{description}
  \item[\textsc{Edit}] (pencil icon). Opens the single-lesson
    edit modal: teacher, subject, classroom, class, optional
    co-teacher. Persists via \texttt{PATCH /api/lessons/\{id\}}.
  \item[\textsc{Move}] (drag onto the destination cell). Calls
    \texttt{POST /api/lessons/\{id\}/move} with body
    \texttt{\{day\_idx, hour\_idx, on\_conflict\}}; the backend
    returns \texttt{200} when no conflict, \texttt{409}
    otherwise.
  \item[\textsc{Unbind}] (arrow icon). Removes the lesson from
    the grid but does not delete it from the database: the
    lesson lands in the pool. Calls \texttt{POST
    /api/lessons/\{id\}/unschedule}.
  \item[\textsc{Delete}] (trash icon, confirmation required).
    Removes the \texttt{Lesson} from the database. Endpoint
    \texttt{DELETE /api/lessons/\{id\}}.
\end{description}

\section{The conflict modal: \emph{Replace or Cancel}}

\subsection{When it fires}

The user drags a lesson (or a pool item) onto a slot that is
already \emph{busy} from the backend's perspective: teacher
busy, class busy, room busy. The frontend first attempts a
\emph{dry-run} (\texttt{POST
/api/lessons/\{id\}/move?on\_conflict=dry\_run}); the backend
answers \texttt{409} with a payload of the form
\begin{lstlisting}[basicstyle=\ttfamily\footnotesize]
{
  "detail": {
    "conflicts": {
      "teacher_busy":  [ ... ],
      "class_busy":    [ ... ],
      "room_busy":     [ ... ]
    }
  }
}
\end{lstlisting}
At this point the \texttt{ScheduleConflictModal} opens (commit
\texttt{59ba4a8}).

\subsection{The three responses, and why drop shows only two}

The modal is canonical: anywhere the application detects a hard
conflict, it shows the same component. Three responses are
possible:
\begin{itemize}
  \item \textsc{Cancel}: closes the modal without action; the
    lesson reverts to its starting position.
  \item \textsc{Unbind}: removes from the calendar \emph{only}
    the conflicting resources --- for room conflicts, clears
    \texttt{classroom\_name}; for teacher or class, deletes the
    conflicting lesson. The lesson being dropped settles where
    intended.
  \item \textsc{Delete} / \textsc{Replace}: deletes the
    conflicting \texttt{Lesson} rows entirely and places the
    new lesson in their stead.
\end{itemize}
In the \emph{drop on occupied slot} flow of \texttt{/schedule}
the \textsc{unbind} option is hidden: placing a new lesson
\emph{without} unbinding the others necessarily leads to a full
replacement, and the dichotomy \emph{Replace / Cancel} best
describes the action (prop \texttt{showUnbind=false},
\texttt{deleteLabel="Sostituisci"}).

\section{The working-hours tab (\texttt{/ore})}

\grokfig{shot_ore}{The working-hours tab: here you define the school
week (days and time bands). Edits propagate to the calendar in real
time through \texttt{workingHoursStore}.}
\label{sec:tab-ore-en}

\subsection{Why a dedicated tab}

Italian school timetables do not share a universal hourly grid:
some schools run a short week (Monday to Friday), others
include Saturday; some use 60-minute periods, others 55 or 50;
some allow variable-length breaks. Hard-coding \texttt{DAYS =
['Mon', \dots, 'Sat']} and \texttt{HOURS = 6}, as the engine
did until March MMXXVI, meant pretending the world conformed
to one example.

Three refactor commits changed the picture: \texttt{b24fdbd}
(data model: \texttt{WorkingDay} and \texttt{WorkingHourSlot}
tables), \texttt{a3c12fb} (REST router), and \texttt{bf6cfa2}
(engine: read configuration from the new tables, drop the
\texttt{DAYS} and \texttt{HOURS} constants). On top of that,
commit \texttt{b3236ab} added the working-hours tab as a
first-class navbar entry.

\subsection{Data model and API}

\begin{itemize}
  \item \texttt{WorkingDay(id, dow, name, active)}: the
    school-active days of the week, each with a readable
    \emph{name} (`Monday', `Wednesday', \dots).
  \item \texttt{WorkingHourSlot(id, day\_id, idx, label,
    start\_min, end\_min, kind)}: the slot inside a day, with
    label (1\textsuperscript{st} hour, 2\textsuperscript{nd}
    hour, recess, \dots), start and end in minutes since
    midnight, and \emph{kind} in
    \textsc{lesson}/\textsc{recess}/\textsc{break}.
\end{itemize}

The REST router \texttt{/api/working-hours/*} exposes:
\begin{lstlisting}[basicstyle=\ttfamily\footnotesize]
GET    /api/working-hours/config
GET    /api/working-hours/days
POST   /api/working-hours/days
PUT    /api/working-hours/days/{day_id}
DELETE /api/working-hours/days/{day_id}            # cascade slots
PUT    /api/working-hours/days/{day_id}/slots      # replace
POST   /api/working-hours/reset                    # default Mon-Sat 8-14
\end{lstlisting}

The \texttt{/reset} endpoint restores the default configuration:
\emph{Monday-Saturday, six consecutive 60-minute periods, from
08:00 to 14:00}.

\subsection{The vintage-style UI}

Commits \texttt{d11e396} and \texttt{d4bc873} turned the slot
editor into a typographically pleasing worksheet, freely inspired
by early-20th-century handwritten school registers: sober serif
type, no baroque-coloured buttons, a single accent of colour
(indigo) for actions, three distinct CSS cursors to carry the
drag behaviour without explaining it in words.

\begin{itemize}
  \item Each existing-slot cell shows two restrained buttons:
    \emph{Edit} (opens an inline popover that floats
    \emph{above} the cell, never below, to keep the reading
    flow clean) and \emph{Delete} (with a double-click to
    confirm).
  \item \emph{Drop key hints} guide the user: dragging on the
    bottom edge of the cell \emph{resizes} the slot, dragging
    from its body \emph{moves} it. Hints appear only during a
    drag.
  \item The \emph{translucent ghost}: when a new slot is being
    drawn, a semi-transparent preview follows the cursor.
    This is what commit \texttt{d4bc873} calls the
    \emph{ghost element}.
\end{itemize}

\subsection{The fix that unblocked the Cypress suite}

A side note, but instructive. Commit \texttt{5c1ba34} carries
the message ``\emph{namespace-import was masking api.get}''.
Translation: the \texttt{ore/+page.svelte} file imported
\texttt{api} as a namespace (\texttt{import * as api}), but a
local \texttt{api} variable declared further down shadowed the
namespace, breaking \texttt{api.get(\dots)} at runtime. The fix
(\texttt{import \{ api \} from \dots}) carried the working-hours
Cypress suite from \texttt{0/15} to \texttt{15/15} green. The
lesson: namespace imports are fragile in the presence of
homonymous variables; named imports are safer.

\section{The shared store \texttt{workingHoursStore}}
\label{sec:store-en}

Until commit \texttt{e0b76bc}, changing a slot in \texttt{/ore}
did not propagate to the other calendar views until the page
reloaded. Typical example: I add the 7\textsuperscript{th}
hour in the working-hours tab, switch to \texttt{/schedule},
and the grid still shows six columns.

The fix is a Writable Svelte store, \texttt{workingHoursStore},
that holds the bundle of \texttt{WorkingDay} and
\texttt{WorkingHourSlot} and auto-updates every time the
working-hours tab confirms a save. Every
\texttt{WeeklyCalendarView} subscribed to the store receives
the new slot array and redraws. No reload, no prop drilling: the
propagation is \emph{live}.

\section{Bulk operations on \texttt{/assignments}}
\label{sec:bulk-en}

Commits \texttt{e6314eb} (backend) and \texttt{16f882a}
(frontend) completed the assignments page with a \emph{bulk
operations} panel. Multi-select existed elsewhere; here it
unfolds into five actions:

\begin{itemize}
  \item \textsc{Delete}: removes the selected rows
    (\texttt{POST /api/assignments/bulk-delete}).
  \item \textsc{Lock}/\textsc{Unlock}: raises or lowers the
    \texttt{locked} flag (\texttt{POST
    /api/assignments/bulk-set-flag} with \texttt{flag=locked}).
  \item \textsc{Change teacher}: bulk-reassigns the teacher
    (\texttt{POST /api/assignments/bulk-change-teacher}).
  \item \textsc{Set flag}: a generic version of \emph{lock} for
    arbitrary boolean flags supported by the schema.
\end{itemize}

The user selects rows (Ctrl-click, Shift-click, or the header
checkbox for ``select all''), opens the panel, confirms; the
frontend issues \emph{a single} request to the backend, which
opens a transaction and applies the operation in bulk. For
schools with hundreds of assignments --- common at upper
secondary level --- the difference compared with hundreds of
serial calls is one order of magnitude.

\section{The whole picture}

The four innovations described in this chapter are not
independent: they merge into a single editorial experience.

\fleuron

The working-hours tab configures the grid. The shared store
propagates it. \texttt{WeeklyCalendarView} renders it. The
conflict modal resolves the last ambiguity when the user
composes, dismantles and recomposes the school week with a
finger. In the small things of the interaction --- an
\texttt{ns-resize} cursor on the edge, a red preview when the
drop would be wrong, a fleuron in place of a square bullet ---
lies the difference between a web page and a school manual.
